\documentclass[12pt]{article}
\usepackage[utf8]{inputenc}
\usepackage{amsmath}
\usepackage{amsfonts}
\usepackage{amssymb}
\usepackage[margin=2cm]{geometry}
\geometry{a4paper, margin=1in}

\begin{document}

\begin{center}
    Universidad de Antioquia - Instituto de Física\\
    Mecánica de Medios Continuos\\
    Primer examen (23\%) 24/03/2026
\end{center}

\textbf{NOTA:} es necesario hacer todos los procedimientos matemáticos y las respuestas deben ser simplificadas y justificadas.

\begin{enumerate}
    \item (40\%) Considere una atmósfera en equilibrio hidrostático sometida a un campo gravitacional constante $\vec{g} = -g_0 \hat{k}$. La atmósfera está compuesta por dos capas de gas ideal:
    \begin{itemize}
        \item para $0 \le z \le z_0$, la temperatura es constante e igual a $T_1$;
        \item para $z > z_0$, la temperatura es constante e igual a $T_2$.
    \end{itemize}
    Suponga que el gas tiene la misma composición química en ambas capas, de modo que la constante específica del gas es la misma en toda la atmósfera. Suponga además que la presión es continua en la interfaz $z = z_0$.
    \begin{itemize}
        \item[a)] (15\%) Obtenga la presión $p(z)$ en cada una de las dos regiones.
        \item[b)] (15\%) Obtenga la densidad $\rho(z)$ en cada región.
        \item[c)] (5\%) Defina la altura de escala en cada capa e interprétela físicamente.
        \item[d)] (5\%) Determine el cociente $\frac{\rho(z_0^+)}{\rho(z_0^-)}$ y explique su significado físico.
    \end{itemize}

    \item (30\%) Una capa infinita de fluido incompresible y densidad constante $\rho$ ocupa la región $|z| < a$, extendiéndose infinitamente en las direcciones $x$ e $y$. Debido a la simetría traslacional en los ejes $x$ e $y$, el problema es unidimensional y todas las cantidades físicas dependen únicamente de la coordenada $z$. Determine la distribución de presión en el interior del fluido y demuestre que la presión central es
    \[
    p(0) = 2\pi G \rho^2 a^2.
    \]

    \item (30\%) Considere un fluido politrópico descrito por la ecuación de estado
    \[
    p = \alpha \rho^n,
    \]
    en reposo en un sistema de referencia que rota con velocidad angular constante $\Omega$ alrededor del eje $z$, y en presencia de un campo gravitacional uniforme
    \[
    \vec{g} = -g_0 \hat{k}.
    \]
    Trabaje en coordenadas cilíndricas.
    \begin{itemize}
        \item[a)] (20\%) Use la ecuación de estado para obtener una ecuación diferencial para $\rho(r, z)$ y determine la distribución de densidad, sabiendo que
        \[
        \rho(0, 0) = \rho_0.
        \]
        \item[b)] (10\%) Discuta la región del espacio donde la solución es físicamente aceptable. ¿Cómo depende esta región de los parámetros del problema?
    \end{itemize}
\end{enumerate}

\end{document}
